\documentclass[main.tex]{subfiles}

\begin{document}

\chapter*{Conclusioni}
  \addcontentsline{toc}{chapter}{Conclusioni}
  Lo sviluppo dell'infrastruttura software
  per il rivelatore SAND necessitava di strumenti in grado di testare le
  pipeline di analisi ancor prima che i complessi algoritmi di ricostruzione
  delle tracce fossero ultimati.
  Inoltre, i file di output \textit{Common Analysis File} (CAF)
  rappresentano delle strutture dati nuove per la
  collaborazione e in via di sviluppo, che richiedevano un meccanismo di
  popolamento affidabile per poterne studiare a fondo il design e
  individuarne eventuali storture.

  Per risolvere queste criticità, è stato sviluppato il modulo
  fake-reco all'interno del nuovo microFrameWork (ufw):
  un framework \textit{object-oriented} sviluppato ad hoc nell'ultimo anno
  dalla sezione INFN di Bologna.
  Insieme ad esso, sono state implementate interfacce specifiche
  (\texttt{genie\_reader}, \texttt{edep\_reader}, \texttt{root\_tgeomanager})
  per leggere la geometria del rivelatore e i formati di output dei generatori.
  Il software creato funge da ponte: traduce le strutture dati di EDep-Sim e
  GENIE copiando le informazioni della verità Monte Carlo direttamente negli
  oggetti di ricostruzione finali previsti dai file CAF\@.

  Grazie a questo lavoro, è ora disponibile un workflow operativo in grado
  di generare file CAF strutturalmente completi.
  Sebbene i dati nei campi di ricostruzione siano attualmente ``perfetti'' (in
  quanto copie esatte della simulazione), essi offrono una base fondamentale per
  permettere ai ricercatori di familiarizzare con i nuovi formati e per testare
  a fondo i software di analisi dati.

  La roadmap prevede l'evoluzione del modulo attuale in fast-reco,
  introducendo un opportuno \textit{smearing} ai dati di ricostruzione per
  simulare in modo più realistico la risposta del rivelatore e le incertezze di
  misura.
  Come passo successivo verso la ricostruzione completa, il modulo verrà poi
  frammentato in processi più piccoli, dedicati ai singoli sottorivelatori di
  SAND (ECAL, STT, GRAIN); questo accorgimento permetterà di sostituire
  gradualmente i vari blocchi di fast-reco con i reali algoritmi di
  ricostruzione non appena questi saranno disponibili nel nuovo framework,
  mantenendo l'intera catena di analisi costantemente funzionante.

\end{document}